\section{Cumulative distribution functions}

The cumulative distribution function is the most universal way to describe a
real-valued random variable. It works for discrete, continuous, and mixed
probability laws.

\begin{definitionbox}
For a real-valued random variable \(X\), the \textbf{cumulative distribution
function}, or \textbf{CDF}, is
\[
F_X(x)=P(X\leq x),
\qquad x\in\mathbb R.
\]
\end{definitionbox}

The value \(F_X(x)\) is the probability accumulated at or below \(x\).

\subsection*{CDF from a density}

If \(X\) has density \(f_X\), then
\[
F_X(x)=\int_{-\infty}^{x}f_X(t)\,dt.
\]
The variable \(t\) is a dummy integration variable. The CDF is therefore an
accumulated area function.

\subsection*{Recovering interval probabilities}

For \(a<b\),
\[
\{X\leq b\}=\{X\leq a\}\cup\{a<X\leq b\},
\]
and the two events on the right are disjoint. Hence
\[
P(a<X\leq b)=F_X(b)-F_X(a).
\]
If \(X\) has a density, then point probabilities are zero, so endpoint choices do
not change the result:
\[
P(a\leq X\leq b)=F_X(b)-F_X(a).
\]

\begin{theorembox}
For a random variable with a density,
\[
F_X(x)=\int_{-\infty}^{x}f_X(t)\,dt,
\qquad
P(a\leq X\leq b)=F_X(b)-F_X(a).
\]
\end{theorembox}

\subsection*{Example: uniform distribution on \([0,1]\)}

For \(X\sim\operatorname{Unif}(0,1)\),
\[
f_X(x)=
\begin{cases}
1, & 0\leq x\leq1,\\
0, & \text{otherwise}.
\end{cases}
\]
Therefore
\[
F_X(x)=
\begin{cases}
0, & x<0,\\
x, & 0\leq x\leq1,\\
1, & x>1.
\end{cases}
\]
For example,
\[
P(0.2\leq X\leq0.7)=F_X(0.7)-F_X(0.2)=0.5.
\]

\subsection*{Example: a triangular density}

Let
\[
f_X(x)=
\begin{cases}
2x, & 0\leq x\leq1,\\
0, & \text{otherwise}.
\end{cases}
\]
Then
\[
F_X(x)=
\begin{cases}
0, & x<0,\\
x^2, & 0\leq x\leq1,\\
1, & x>1.
\end{cases}
\]
Consequently,
\[
P(0.5\leq X\leq1)=1-0.5^2=0.75.
\]

\subsection*{Basic properties of every CDF}

Every cumulative distribution function satisfies:
\begin{itemize}
    \item \(0\leq F_X(x)\leq1\) for every \(x\);
    \item \(F_X\) is non-decreasing;
    \item \(F_X\) is right-continuous;
    \item \(\lim_{x\to-\infty}F_X(x)=0\);
    \item \(\lim_{x\to\infty}F_X(x)=1\).
\end{itemize}

Right-continuity is a property of every CDF, not only of density models. When
\(X\) has a density, its CDF is absolutely continuous and therefore continuous.
At every point where \(f_X\) is continuous, the fundamental theorem of calculus
gives
\[
F_X'(x)=f_X(x).
\]

\begin{remarkbox}
The density is the local rate at which the CDF accumulates probability. A high
density makes the CDF rise quickly; a low density makes it rise slowly.
\end{remarkbox}

\subsection*{CDF versus density}

The density and the CDF describe the same distribution from different points of
view. The density describes local concentration per unit length. The CDF records
accumulated probability. A density need not exist for every probability law, but
a CDF always exists.

\begin{summarybox}
The CDF is the universal description
\[
F_X(x)=P(X\leq x).
\]
For a density model it is obtained by integration, while interval probabilities
are recovered by differences of CDF values.
\end{summarybox}
